
% Montaje experimental para la caracterización angular.
% El entorno figure, caption y label se definen en el capítulo.

\begin{tikzpicture}[
    >=Stealth,
    line cap=round,
    line join=round,
    font=\small,
    x=1cm,
    y=1cm
]

    % =================================================
    % 1. GEOMETRIA DEL ARREGLO
    % =================================================

    % Centro del arreglo desplazado hacia la izquierda
    \coordinate (O) at (-3,0);

    % Micrófonos: mayor separación horizontal
    \coordinate (Mone) at (-3.425,0);
    \coordinate (Mtwo) at (-2.575,0);

    % Eje horizontal principal
    \draw[line width=0.95pt]
        (-5.35,0) -- (-0.65,0);

    % Eje vertical principal
    \draw[line width=0.95pt]
        (-3,-0.05) -- (-3,4.35);

    % =================================================
    % 2. POSICIONES ANGULARES
    % =================================================

    \foreach \ang in {-60,-45,-30,-15,15,30,45,60}{
        \draw[
            densely dashed,
            line width=0.8pt
        ]
        ({-3+4.0*sin(\ang)},{4.0*cos(\ang)})
        --
        ({-3+0.38*sin(\ang)},{0.38*cos(\ang)});
    }

    % =================================================
    % 3. FUENTE ACUSTICA
    % =================================================

    % Posición ilustrativa de la fuente a +30 grados
    \coordinate (S) at (
        {-3+3.55*sin(30)},
        {3.55*cos(30)}
    );

    \node[
        circle,
        draw=orange!85!black,
        fill=orange!18,
        line width=1pt,
        minimum size=6mm,
        inner sep=0pt
    ] at (S) {};

    \node[
        anchor=south,
        inner sep=2pt
    ] at (-1.22,3.48) {Fuente};

    % =================================================
    % 4. REFERENCIA ANGULAR
    % =================================================

    \draw[
        line width=0.85pt
    ]
    (-3,0.78)
    arc[
        start angle=90,
        end angle=60,
        radius=0.78
    ];

    \node[
        fill=white,
        inner sep=2pt
    ] at (-2.77,0.83) {$\theta$};

    % =================================================
    % 5. DISTANCIA RADIAL
    % =================================================

    \node[
        fill=white,
        inner sep=2pt
    ] at (-4.05,2.35)
    {$r=30\,\mathrm{cm}$};

    % =================================================
    % 6. MICROFONOS
    % =================================================

    % Micrófono 1
    \node[
        circle,
        draw=black!45,
        fill=black!12,
        line width=0.9pt,
        minimum size=5.8mm,
        inner sep=0pt
    ] at (Mone) {};

    % Micrófono 2
    \node[
        circle,
        draw=black!45,
        fill=black!12,
        line width=0.9pt,
        minimum size=5.8mm,
        inner sep=0pt
    ] at (Mtwo) {};

    % Identificación individual con líneas guía
    \draw[black!65,line width=0.55pt]
        (-3.425,-0.15) -- (-3.72,-0.48);

    \draw[black!65,line width=0.55pt]
        (-2.575,-0.15) -- (-2.28,-0.48);

    \node[
        anchor=north
    ] at (-3.76,-0.48) {Mic. 1};

    \node[
        anchor=north
    ] at (-2.24,-0.48) {Mic. 2};

    % =================================================
    % 7. DISTANCIA ENTRE MICROFONOS
    % =================================================

    \draw[densely dashed,black!45]
        (-3.425,-0.22) -- (-3.425,-1.12);

    \draw[densely dashed,black!45]
        (-2.575,-0.22) -- (-2.575,-1.12);

    % Cota de separación
    \draw[
        <->,
        line width=0.75pt
    ]
        (-3.425,-1.02) -- (-2.575,-1.02);

    \node[
        fill=white,
        inner sep=2pt
    ] at (-3,-1.30)
    {$d=5\,\mathrm{cm}$};

    % Centro geométrico
    \fill (O) circle (1.4pt);

    % =================================================
    % 8. CONFIGURACION EXPERIMENTAL
    % =================================================

    \node[
        draw=black!65,
        rounded corners=2pt,
        fill=black!2,
        line width=0.65pt,
        text width=4.15cm,
        align=left,
        inner sep=9pt,
        anchor=west
    ] at (1.05,1.35) {%
        \textbf{Configuración experimental}\\[6pt]
        Línea base: $5\,\mathrm{cm}$\\[3pt]
        Distancia nominal: $30\,\mathrm{cm}$\\[3pt]
        Ángulos: $-60^\circ$ a $+60^\circ$\\[3pt]
        Paso angular: $15^\circ$
    };

\end{tikzpicture}
